\clearpage
\VerdierPhysicalPageAnchor{0080}%
\VerdierSemanticAnchor{ch1-2-5-10}
\noindent\textbf{2.5.10.}\quad
\VerdierSemanticAnchor{ch1-2-5-10-multiadditive-setup}
设
$F\colon\prod_{i\in[n]}C_i\longrightarrow C$ 是一个多加法函子。记

\VerdierSemanticAnchor{ch1-2-5-10-quotient-functors}
\begin{equation*}
  Q_i\colon\operatorname{comp}(C_i)\longrightarrow K(C_i)
\end{equation*}
为典范函子。由 (2.5.6) 与 (2.5.7) 可知，当 $C$ 中的可数乘积可表示时，
存在唯一的函子

\VerdierSemanticAnchor{ch1-2-5-10-pi-extension}
\begin{equation*}
  K_{\Pi}(F)\colon
  \prod_{i\in[n]}K(C_i)\longrightarrow K(C)
\end{equation*}
使下图交换：

\VerdierSemanticAnchor{ch1-2-5-10-eq1}
\begin{equation}
  \resizebox{0.84\textwidth}{!}{%
    $\begin{tikzcd}[ampersand replacement=\&,column sep=6.8em,row sep=5.0em]
      \displaystyle\prod_{i\in[n]}\operatorname{comp}(C_i)
        \arrow[r,"{\operatorname{comp}_{\Pi}F}"]
        \arrow[d,"{\prod_{i\in[n]}Q_i}"']
        \& \operatorname{comp}(C)
          \arrow[d,"Q"] \\
      \displaystyle\prod_{i\in[n]}K(C_i)
        \arrow[r,"{K_{\Pi}F}"]
        \& K(C)\mathpunct{.}
    \end{tikzcd}$%
  }
  \tag{2.5.10.1}
\end{equation}

\VerdierSemanticAnchor{ch1-2-5-10-twisted-structure}
此外，由于函子 $\prod_{i\in[n]}Q_i$ 是严格 $\mathbb Z[n]$-函子，
函子 $K_{\Pi}F$ 上存在唯一的、由 $\mathbb Z[n]$ 的有序基
$e_1,\ldots,e_n$ 所对应的 Koszul 余循环扭曲的分次
$\mathbb Z[n]$-函子结构，使下列扭曲 $\mathbb Z[n]$-函子之间的等式成立：

\VerdierSemanticAnchor{ch1-2-5-10-eq2}
\begin{equation}
  K_{\Pi}F\prod_{i\in[n]}Q_i
  =Q\operatorname{comp}_{\Pi}F.
  \tag{2.5.10.2}
\end{equation}

\VerdierSemanticAnchor{ch1-2-5-10-morphism-setup}
最后，设 $F$ 与 $F'$ 是两个多加法函子，
$m\colon F\longrightarrow F'$ 是一个函子态射。由 (2.5.7) 可知，
存在唯一的函子态射：

\VerdierSemanticAnchor{ch1-2-5-10-eq3}
\begin{equation}
  K_{\Pi}m\colon K_{\Pi}F\longrightarrow K_{\Pi}F'
  \tag{2.5.10.3}
\end{equation}
使得：

\VerdierSemanticAnchor{ch1-2-5-10-eq4}
\begin{equation}
  K_{\Pi}m\star\prod_{i\in[n]}Q_i
  =Q\star\operatorname{comp}_{\Pi}m.
  \tag{2.5.10.4}
\end{equation}

\VerdierSemanticAnchor{ch1-2-5-10-pi-terminology}
函子 $K_{\Pi}F$ 称为函子 $F$ 在复形同伦范畴上的
$\Pi$-\emph{扩张}。态射 $K_{\Pi}m$ 称为态射 $m$
在复形同伦范畴上的 $\Pi$-\emph{扩张}。

\VerdierSemanticAnchor{ch1-2-5-10-sigma-analogue}
当 $C$ 中的可数直和可表示时，以类似方式定义函子 $K_{\Sigma}F$，
即函子 $F$ 在复形同伦范畴上的 $\Sigma$-\emph{扩张}；以及态射
$K_{\Sigma}m$，即态射 $m$ 在复形同伦范畴上的
$\Sigma$-\emph{扩张}。
